% 数学ノート集・共通テンプレート v2（積分論2026年度第1部の版面を参考）
% A5判。章内で定義・定理・例・注意等の番号を共通化。
% コンパイル： platex -interaction=nonstopmode math_note_template_v2.tex (2回)
%              dvipdfmx math_note_template_v2.dvi
\documentclass[a5paper,10pt,onecolumn]{jsarticle}
\usepackage[dvipdfmx,top=22mm,bottom=20mm,inner=19mm,outer=18mm,headheight=13pt,headsep=6mm]{geometry}
\usepackage{amsmath,amssymb,amsthm,mathtools}
\usepackage{enumitem,fancyhdr,titlesec}
\usepackage[dvipdfmx,bookmarks=true,colorlinks=false]{hyperref}
\usepackage{pxjahyper}

% ----- 差し替える文書情報 -----
\newcommand{\NoteTitle}{位相群の基本構造}
\newcommand{\NoteSeries}{数学ノート集}

% ----- 文字・組版 -----
% 参考PDFの質素な明朝＋Computer Modernの印象を引き継ぎながら，
% A5判で窮屈にならないよう本文の行間を少し確保する．
\linespread{1.13}
\setlength{\parindent}{1zw}
\setlength{\parskip}{0pt}
\setlength{\abovedisplayskip}{6pt plus 2pt minus 2pt}
\setlength{\belowdisplayskip}{6pt plus 2pt minus 2pt}
\setlength{\abovedisplayshortskip}{3pt plus 2pt minus 1pt}
\setlength{\belowdisplayshortskip}{4pt plus 2pt minus 1pt}
\setlist[enumerate]{topsep=3pt,itemsep=1pt,parsep=0pt,leftmargin=2.4zw}
\titleformat{\section}{\normalfont\normalsize\bfseries}{\thesection}{0.8zw}{}
\titlespacing*{\section}{0pt}{2.0ex plus .4ex}{.9ex}
\titleformat{\subsection}{\normalfont\normalsize\bfseries}{\thesubsection}{0.7zw}{}
\titlespacing*{\subsection}{0pt}{1.5ex plus .3ex}{.6ex}
\raggedbottom

% ----- 柱とページ番号：ページ上部、罫線なし -----
\fancypagestyle{mathnote}{%
  \fancyhf{}%
  \fancyhead[L]{\footnotesize \NoteSeries\ （\NoteTitle）}%
  \fancyhead[R]{\footnotesize \thepage}%
  \renewcommand{\headrulewidth}{0pt}%
  \renewcommand{\footrulewidth}{0pt}%
}
\pagestyle{mathnote}

% ----- 節内通し番号の定義・定理群 -----
\newtheoremstyle{mathnotetheorem}
 {4.5pt}{4.5pt}{\normalfont}{}%
 {\bfseries}{}{0.65zw}%
 {\thmname{#1}\thmnumber{ #2}\thmnote{（#3）}}
\theoremstyle{mathnotetheorem}
\newtheorem{theorem}{定理}[section]
\newtheorem{definition}[theorem]{定義}
\newtheorem{lemma}[theorem]{補題}
\newtheorem{proposition}[theorem]{命題}
\newtheorem{corollary}[theorem]{系}
\newtheorem{example}[theorem]{例}
\newtheorem{remark}[theorem]{注意}
\renewcommand{\proofname}{証明}
\renewcommand{\qedsymbol}{$\square$}
\makeatletter
\renewenvironment{proof}[1][\proofname]{%
  \par\pushQED{\qed}\normalfont\topsep4\p@\@plus2\p@\relax
  \trivlist\item[\hskip\labelsep\bfseries #1]\ignorespaces
}{\popQED\endtrivlist\@endpefalse}
\makeatother

% ----- 使うノートで適宜拡張する記号 -----
\newcommand{\R}{\mathbb{R}}
\newcommand{\norm}[1]{\lVert #1\rVert}
\newcommand{\abs}[1]{\lvert #1\rvert}

\begin{document}
\section*{読了後にできること}
本稿を読むための前提は，位相群，積位相，部分群と商群．
本稿の具体的な到達課題は，具体的な群の位相群構造を証明し，閉部分群と商位相を確認する．


\section{位相群の定義}

本稿では，特に断らない限り位相空間をHausdorffとする．
群 $G$ の単位元を $e$ と書く．積 $G\times G$ には積位相を入れる．

\begin{definition}
群 $G$ に位相が与えられているとき，積写像
\[
 m:G\times G\longrightarrow G,\qquad m(x,y)=xy
\]
および逆元写像
\[
 \iota:G\longrightarrow G,\qquad \iota(x)=x^{-1}
\]
が連続であれば，$G$ を\emph{位相群}という．
\end{definition}

\begin{proposition}[演算の連続性の判定]\label{prop:single}
群 $G$ が位相群であるための必要十分条件は
\[
 q:G\times G\longrightarrow G,\qquad
 q(x,y)=xy^{-1}
\]
が連続であることである．
\end{proposition}
\begin{proof}
$G$ が位相群なら $q=m\circ(\operatorname{id}\times\iota)$
だから連続である．
逆に $q$ が連続なら
$\iota(y)=q(e,y)$ は連続である．
従って $m(x,y)=q(x,\iota(y))$ も連続である．
\end{proof}

\begin{example}
$\mathbb R^n$ は加法と通常の位相により位相群である．
$GL_n(\mathbb R)$ は行列積とEuclid位相の相対位相により
位相群である．
\end{example}
\begin{proof}
$\mathbb R^n$ の加法と符号反転は連続である．
行列積の各成分は行列成分の多項式なので連続である．
$GL_n(\mathbb R)=\{A:\det A\ne0\}$ は開集合であり，
逆行列の成分は余因子を $\det A$ で割った関数なので
$GL_n(\mathbb R)$ 上連続である．
\end{proof}

\begin{lemma}[単位元の近傍]\label{lem:nbhd}
位相群 $G$ の単位元 $e$ の任意の開近傍 $U$ に対し，
次を満たす対称な開近傍 $V$ が存在する：
\[
 e\in V,\qquad V^{-1}=V,\qquad VV\subset U.
\]
さらに，$U$ に含まれる対称な開近傍も存在する．
\end{lemma}
\begin{proof}
$m(e,e)=e$ と積写像の連続性から，
$e$ の開近傍 $W_1,W_2$ が存在して
$W_1W_2\subset U$ となる．
$W=W_1\cap W_2$ と置けば $WW\subset U$ である．
逆元写像は同相写像なので $W^{-1}$ は開である．
\[
 V=W\cap W^{-1}
\]
は対称な開近傍であり $VV\subset WW\subset U$ となる．
最後の主張は $U\cap U^{-1}$ を取ればよい．
\end{proof}

\section{平行移動と連続準同型}

\begin{proposition}[平行移動と共役]\label{prop:translations}
$G$ を位相群とし，$a\in G$ を固定する．
左平行移動 $L_a(x)=ax$，右平行移動 $R_a(x)=xa$，
共役写像 $c_a(x)=axa^{-1}$ はいずれも同相写像である．
\end{proposition}
\begin{proof}
積写像と定値写像の連続性から $L_a,R_a$ は連続であり，
それぞれ逆写像 $L_{a^{-1}},R_{a^{-1}}$ も連続である．
$c_a=L_a\circ R_{a^{-1}}$ も同相写像である．
\end{proof}

\begin{corollary}\label{cor:local}
位相群の任意の点 $a$ の近傍系は，単位元の近傍系を
左平行移動して得られる．
\end{corollary}
\begin{proof}
$L_a$ は $e$ を $a$ に写す同相写像だからである．
\end{proof}

\begin{proposition}[連続性の一点判定]\label{prop:hom-cont}
$G,H$ を位相群，$\varphi:G\to H$ を群準同型とする．
$\varphi$ が単位元 $e_G$ で連続なら全域で連続である．
\end{proposition}
\begin{proof}
任意の $a\in G$ と $a$ に収束するネット $(x_\lambda)$
を取る．
積と逆元の連続性から
$a^{-1}x_\lambda\to e_G$ である．
従って仮定により
$\varphi(a^{-1}x_\lambda)\to e_H$ となる．
\[
 \varphi(x_\lambda)
 =\varphi(a)\varphi(a^{-1}x_\lambda)
 \longrightarrow\varphi(a)
\]
だから $a$ でも連続である．
同じ議論は単位元の近傍を用いても証明できる．
\end{proof}

\begin{proposition}[核]\label{prop:kernel}
$G,H$ をHausdorff位相群，$\varphi:G\to H$ を
連続準同型とする．このとき
$\ker\varphi$ は $G$ の閉正規部分群である．
\end{proposition}
\begin{proof}
群準同型の核が正規部分群であることは群論から従う．
$H$ はHausdorffだから単元集合 $\{e_H\}$ は閉である．
従って
$\ker\varphi=\varphi^{-1}(\{e_H\})$ は閉である．
\end{proof}

\section{部分群の閉包}

\begin{theorem}[部分群の閉包]\label{thm:closure-subgroup}
$H$ を位相群 $G$ の部分群とする．
このとき閉包 $\overline H$ も部分群である．
さらに $H\triangleleft G$ なら
$\overline H\triangleleft G$ である．
\end{theorem}
\begin{proof}
積写像の連続性によって
\[
 \overline H\,\overline H
 \subset\overline{HH}=\overline H
\]
が成り立つ．
実際 $x,y\in\overline H$ と $xy$ の任意の近傍 $U$ を取ると，
積の連続性によって
$x,y$ の近傍 $V,W$ で $VW\subset U$ となる．
$V\cap H$，$W\cap H$ は非空だから
$U\cap HH$ も非空となり，$xy\in\overline{HH}$ である．
また逆元写像は同相なので
$(\overline H)^{-1}=\overline{H^{-1}}=\overline H$ である．
単位元も $\overline H$ に属するから部分群である．

$H$ が正規なら任意の $g\in G$ について
$gHg^{-1}=H$ であり，
共役写像は同相なので
\[
 g\overline H g^{-1}
 =\overline{gHg^{-1}}
 =\overline H.
\]
従って $\overline H$ は正規部分群である．
\end{proof}

\begin{lemma}[閉近傍の構成]\label{lem:regularity}
$G$ をHausdorff位相群とする．
単位元の任意の開近傍 $U$ に対し，開近傍 $V$ で
\[
 e\in V\subset\overline V\subset U
\]
となるものが存在する．
\end{lemma}
\begin{proof}
補題\ref{lem:nbhd}から，対称な開近傍 $W$ で
$WW\subset U$ となるものを取る．
$x\in\overline W$ を任意に取る．
$xW$ は $x$ の開近傍だから $xW\cap W\ne\varnothing$ である．
よって $xw_1=w_2$ となる $w_1,w_2\in W$ が存在し，
\[
 x=w_2w_1^{-1}\in WW\subset U.
\]
従って $V=W$ と取ればよい．
\end{proof}

\begin{corollary}[閉部分群による分離]\label{cor:closed-subgroup}
$G$ をHausdorff位相群，$H$ を閉部分群とする．
任意の $x\notin H$ に対し，$x$ を含み $H$ と交わらない
開集合が存在する．
\end{corollary}
\begin{proof}
$G\setminus H$ は開集合であり $x$ を含むから，
そのまま取ればよい．
\end{proof}

\section{商位相群と普遍性}

\begin{definition}
$N\triangleleft G$ を正規部分群とする．
商群 $G/N$ には，標準射影
\[
 \pi:G\longrightarrow G/N
\]
に関する商位相を入れる．
\end{definition}

\begin{lemma}[商射影は開写像]\label{lem:quotient-open}
$U\subset G$ が開なら $\pi(U)$ は $G/N$ の開集合である．
\end{lemma}
\begin{proof}
\[
 \pi^{-1}(\pi(U))=UN=\bigcup_{n\in N}Un
\]
である．各 $Un$ は右平行移動による開集合なので
$UN$ は開である．
商位相の定義から $\pi(U)$ も開となる．
\end{proof}

\begin{theorem}[商位相群]\label{thm:quotient}
$G$ を位相群，$N\triangleleft G$ を正規部分群とする．
このとき商位相を持つ群 $G/N$ の積と逆元は連続である．
さらに $G$ がHausdorffであるとき，
$G/N$ がHausdorffであるための必要十分条件は
$N$ が $G$ の閉集合であることである．
\end{theorem}
\begin{proof}
$\pi$ は補題\ref{lem:quotient-open}により開かつ全射である．
従って $\pi\times\pi:G\times G\to(G/N)\times(G/N)$
も開かつ全射であり，商写像である．
積写像 $m'$ と逆元写像 $\iota'$ は，
\[
 m'\circ(\pi\times\pi)=\pi\circ m,\qquad
 \iota'\circ\pi=\pi\circ\iota
\]
を満たす．
右辺は連続なので，商写像の普遍性から
$m',\iota'$ は連続である．

$G/N$ がHausdorffなら $\{N\}$ は閉集合であり，
$\pi^{-1}(\{N\})=N$ は閉である．
逆に $N$ が閉なら商位相の定義から
$\{N\}$ は $G/N$ で閉となる．
任意の位相群 $K$ において $\{e_K\}$ が閉なら
$K$ はHausdorffであることを示す．
$x\ne y$ に対し $x^{-1}y\ne e_K$ であるから，
$U=K\setminus\{x^{-1}y\}$ は $e_K$ の開近傍である．
補題\ref{lem:nbhd}から対称な開近傍 $V$ で
$V^{-1}V\subset U$ とできる．
もし $xV\cap yV\ne\varnothing$ なら
$xv_1=yv_2$ となる $v_1,v_2\in V$ が存在し，
$x^{-1}y=v_1v_2^{-1}\in V^{-1}V\subset U$
となって矛盾する．
従って $xV,yV$ は互いに交わらない開近傍であり，
$K$ はHausdorffである．
これを $K=G/N$ に適用して結論を得る．
\end{proof}

\begin{theorem}[商位相群の普遍性]\label{thm:quotient-universal}
$G,H$ を位相群，$N\triangleleft G$ とし，
$\varphi:G\to H$ を連続準同型で
$N\subset\ker\varphi$ を満たすものとする．
このとき一意な連続準同型
$\overline\varphi:G/N\to H$ が存在して
$\varphi=\overline\varphi\circ\pi$ となる．
\end{theorem}
\begin{proof}
群の商の普遍性から
$\overline\varphi(gN)=\varphi(g)$ はwell-definedであり，
唯一の群準同型を定める．
さらに $\overline\varphi\circ\pi=\varphi$ が連続であり，
$\pi$ は商写像なので
$\overline\varphi$ は連続である．
\end{proof}

\begin{proposition}[位相的第一同型定理]\label{prop:first-isomorphism}
$\varphi:G\to H$ を連続かつ開な全射準同型とする．
このとき誘導写像
\[
 G/\ker\varphi\longrightarrow H,\qquad
 g\ker\varphi\longmapsto\varphi(g)
\]
は位相群の同型である．
特に $G$ がコンパクト，$H$ がHausdorffで，
$\varphi$ が連続な全射準同型なら同じ結論が成り立つ．
\end{proposition}
\begin{proof}
商位相群の普遍性によって誘導写像は連続であり，
群論の第一同型定理により全単射である．
$\varphi$ が開なら，商射影も開なので，
誘導写像は開である．従って同相である．
後半では $G/\ker\varphi$ はコンパクトで，
連続全単射による $H$ への像はHausdorffだから
一般位相の定理により同相となる．
\end{proof}

\begin{remark}
単なる連続全射準同型が常に開であるとは限らない．
例えば離散位相の $\mathbb R$ から通常の位相の
$\mathbb R$ への恒等写像は連続全射準同型だが，
開写像ではない．
\end{remark}

\section{連結成分と局所コンパクト性}

\begin{definition}
位相群 $G$ の単位元 $e$ を含む連結成分を $G^0$ と書く．
同様に単位元を含む道連結成分を $G_{\mathrm p}$ と書く．
\end{definition}

\begin{theorem}[単位元の連結成分]\label{thm:identity-component}
$G^0$ は閉正規部分群である．
$G_{\mathrm p}$ は正規部分群であり，
\[
 G_{\mathrm p}\subset G^0.
\]
さらに $G$ が局所道連結なら
$G_{\mathrm p}=G^0$ が成り立つ．
\end{theorem}
\begin{proof}
二つの連結集合の積は連結であり，
連続写像による像も連結である．
従って $G^0G^0$ は $e$ を含む連結集合なので
$G^0G^0\subset G^0$ である．
逆元写像による像 $(G^0)^{-1}$ も
$e$ を含む連結集合だから $G^0$ に含まれる．
よって $G^0$ は部分群である．
また連結集合の閉包は連結なので
$\overline{G^0}$ は $e$ を含む連結集合であり，
連結成分の極大性から $\overline{G^0}=G^0$ となる．
共役は $e$ を固定する同相写像なので
$gG^0g^{-1}=G^0$ であり，正規である．

道連結集合についても積・逆元・共役の像の道連結性から
同じ議論で $G_{\mathrm p}$ が正規部分群となる．
道連結集合は連結なので $G_{\mathrm p}\subset G^0$ である．
局所道連結空間では道連結成分が開集合となる．
その補集合は他の道連結成分の和なので開であり，
道連結成分は閉でもある．
$G^0$ は連結で，互いに素な開閉道連結成分へ分割されるので，
一つの道連結成分に含まれなければならない．
$e$ を含むものは $G_{\mathrm p}$ だから
$G^0\subset G_{\mathrm p}$ を得る．
\end{proof}

\begin{definition}
位相空間 $X$ が\emph{局所コンパクト}であるとは，
任意の点がコンパクトな近傍を持つことをいう．
\end{definition}

\begin{proposition}[単位元での局所コンパクト性]\label{prop:local-compact}
Hausdorff位相群 $G$ は，$e$ がコンパクトな近傍を持つとき，
かつそのときに限り局所コンパクトである．
また $G$ が局所コンパクトなら，
任意の $e$ の開近傍 $U$ に対し，
コンパクトな閉包を持つ対称な開近傍 $V$ で
$\overline V\subset U$ となるものが存在する．
\end{proposition}
\begin{proof}
単位元のコンパクトな近傍 $K$ があれば，
平行移動によって各 $g\in G$ は
コンパクトな近傍 $gK$ を持つので局所コンパクトである．
逆向きは定義から従う．
さらに $e$ のコンパクトな近傍 $K$ を取り，
$e\in W\subset K$ となる開集合 $W$ を選ぶ．
$W\cap U$ は $e$ の開近傍だから，
補題\ref{lem:regularity}によって
開近傍 $V_0$ で $\overline{V_0}\subset W\cap U$ となるものを取れる．
$V=V_0\cap V_0^{-1}$ と置けば対称な開近傍であり，
$\overline V\subset\overline{V_0}\subset U$ となる．
また $\overline V$ は $K$ の閉部分集合なのでコンパクトである．
\end{proof}

\section{連続群作用と軌道空間}

\begin{definition}
位相群 $G$ の位相空間 $X$ への\emph{左作用}とは，
連続写像
\[
 a:G\times X\to X,\qquad (g,x)\mapsto gx
\]
で，$ex=x$ と $(gh)x=g(hx)$ を満たすものをいう．
$x\in X$ に対し
$G_x=\{g\in G:gx=x\}$ を安定化部分群，
$Gx=\{gx:g\in G\}$ を軌道という．
\end{definition}

\begin{proposition}[安定化部分群]\label{prop:stabilizer}
$G$ がHausdorff位相群で $X$ がHausdorffなら，
任意の $x\in X$ の安定化部分群 $G_x$ は閉部分群である．
\end{proposition}
\begin{proof}
$G_x$ が部分群であることは作用の公理から従う．
軌道写像 $a_x:G\to X$，$a_x(g)=gx$ は連続であり，
$X$ がHausdorffなので $\{x\}$ は閉である．
従って $G_x=a_x^{-1}(\{x\})$ は閉である．
\end{proof}

\begin{proposition}[軌道と商群]\label{prop:orbit}
$G$ をHausdorff位相群，$X$ をHausdorff空間とし，
$G$ は $X$ に連続に作用するとする．
任意の $x\in X$ に対して
\[
 \Phi:G/G_x\longrightarrow Gx,\qquad gG_x\longmapsto gx
\]
はwell-definedな連続全単射である．
$G$ がコンパクトなら $\Phi$ は同相写像である．
\end{proposition}
\begin{proof}
$gG_x=hG_x$ と $h^{-1}g\in G_x$ は同値であり，
後者は $gx=hx$ と同値だから
$\Phi$ はwell-definedかつ全単射である．
軌道写像 $a_x$ は
$\Phi\circ\pi=a_x$ を満たすので，
商位相の普遍性から $\Phi$ は連続である．
$G$ がコンパクトなら $G/G_x$ もコンパクトであり，
$Gx$ はHausdorff空間 $X$ の部分空間である．
従ってコンパクト空間からHausdorff空間への
連続全単射 $\Phi$ は同相である．
\end{proof}

\begin{example}[回転群と球面]
$G=SO(3)$ が $S^2$ に通常の行列作用で作用するとする．
北極 $e_3$ の安定化部分群は $SO(2)$ と同型であるから，
\[
 SO(3)/SO(2)\cong S^2
\]
は同相である．
\end{example}
\begin{proof}
$SO(3)$ の各元は直交行列なので球面を保ち，
作用は行列積の連続性により連続である．
任意の単位ベクトルを北極へ移す回転が存在するため，
作用は推移的である．
北極を固定する回転は直交補空間 $e_3^\perp$ 上の
向きを保つ直交変換であり，これは $SO(2)$ に等しい．
$SO(3)$ はEuclid空間中の閉有界集合なのでコンパクトである．
命題\ref{prop:orbit}を適用すれば同相が従う．
\end{proof}

\begin{remark}
局所コンパクト群のHaar測度と，連続表現による調和解析は，
ここで定めた位相群と連続作用の自然な発展である．
一方，位相体は加法位相群と乗法構造の両立を要求する対象である．
\end{remark}

\section{到達課題の完全解答}
\noindent\textbf{円周群を位相群として構成する．}
\begin{proposition}
$S^1=\{z\in\mathbb C:|z|=1\}$ は
複素数の積によりコンパクトHausdorff位相群となる．
\end{proposition}
\begin{proof}
$|zw|=|z||w|=1$，$z^{-1}=\bar z$ だから群である．
複素数の積と共役は連続なので
群演算は連続である．
$S^1$ は $\mathbb R^2$ の有界閉集合なので
コンパクトであり，部分空間としてHausdorffである．
\end{proof}
\begin{proposition}
写像 $\mathbb R/\mathbb Z\to S^1$，
$[t]\mapsto e^{2\pi it}$ は位相群同型である．
\end{proposition}
\begin{proof}
指数法則から群準同型であり，
周期性から核は $\mathbb Z$，
像は $S^1$ 全体である．
商写像の普遍性から連続である．
$\mathbb R/\mathbb Z$ は $[0,1]$ の連続像なので
コンパクトであり，
$S^1$ はHausdorffだから連続全単射は同相となる．
\end{proof}
\noindent\textbf{到達確認．}
$S^1$ の有限部分群
$\mu_n=\{z:z^n=1\}$ が閉であることを示せ．
\par\noindent\textbf{解答．}
連続写像 $z\mapsto z^n$ の閉集合 $\{1\}$
の逆像だから閉である．
\end{document}